% Gesamttest «Dynamik» — Quelle des PDFs (python3 scripts/build-lp-pdf.py dynamik).
% Musterlösung und Raster: bewertungspaket.tex. Alle Werte mit python3 nachgerechnet (04.10.2026).
% Jede Aufgabe fragt in eine andere Richtung als die Aufgaben und Übungen der Seite — rückwärts
% (G4, G5, G6), am gegebenen Diagramm (G2) oder als Begründung (G1) —, HOWTO-leitprogramme §9.
\documentclass[11pt]{article}
\usepackage{../lp-druck}
\renewcommand{\lpname}{Dynamik}
\renewcommand{\lpteilgebiet}{4.2}
\renewcommand{\lpfuss}{Gesamttest · 25 Punkte}
\begin{document}
\lpkopf{Gesamttest}{%
  \vspace{4pt}\noindent\begin{tabularx}{\linewidth}{@{}X X@{}}
  Code (kein Name): \hrulefill & Punkte: \hrulefill\ / 25
  \end{tabularx}}

\begin{lpinfo}{Anleitung}
\textbf{Zeit:} rund 25 Minuten. \textbf{Hilfsmittel:} Taschenrechner und Formelsammlung, in allen Teilen.
Schreib zu jeder Aufgabe zuerst die Formel, dann die Werte \textbf{mit Einheiten}. Punkte gibt es für den Weg,
für das Ergebnis mit Einheit und für Begründungen in eigenen Worten. Rechne mit $g = 9.81\;\text{m/s}^2$; Faden und Rollen sind masselos.
Danach: Lösung scannen oder fotografieren (ohne Namen und Standort) und mit dem \emph{Bewertungspaket} bewerten lassen.
\end{lpinfo}

\lpteil{Teil A · Kraft, Masse und Gesamtkraft}{Kapitel 1 und 2 · 12 Punkte}

\begin{lpaufgabe}{G1}{Begriffe}{4 P}
\begin{enumerate}[label=\alph*),leftmargin=*,itemsep=2pt,topsep=2pt]
\item Woran erkennt man, dass auf einen Körper eine Kraft wirkt?
\item Formuliere das Trägheitsgesetz.
\item Die Kraft auf einen Körper wird verdreifacht, gleichzeitig wird seine Masse halbiert. Um welchen Faktor ändert sich die Beschleunigung? Begründe.
\item Eine Astronautin hat mit Anzug die Masse $75\;\text{kg}$. Was davon ändert sich auf dem Mond ($g = 1.62\;\text{m/s}^2$): ihre Masse oder ihre Gewichtskraft? Wie gross ist ihre Gewichtskraft dort?
\end{enumerate}
\lpplatz{12}
\end{lpaufgabe}

\begin{lpaufgabe}{G2}{Autofahrt im Diagramm}{5 P}
Das $v$-$t$-Diagramm zeigt die Fahrt eines Autos mit $m = 1200\;\text{kg}$ auf gerader Strasse.\par\smallskip
\begin{center}\begin{tikzpicture}
\begin{axis}[width=10cm,height=5.2cm,xmin=0,xmax=18,ymin=0,ymax=17.5,axis lines=left,
  xtick={0,1,...,16},ytick={0,3,...,15},grid=both,grid style={lplinie!70},xticklabel style={font=\small},
  xlabel={$t$ [s]},ylabel={$v$ [m/s]},xlabel style={at={(1,0)},anchor=north east},ylabel style={rotate=-90,at={(0,1)},anchor=south east}]
\addplot[very thick,lpgruen] coordinates {(0,0) (5,15) (12,15) (15,0)};
\end{axis}
\end{tikzpicture}\end{center}
\begin{enumerate}[label=\alph*),leftmargin=*,itemsep=2pt,topsep=2pt]
\item Wie gross ist die Gesamtkraft auf das Auto in den ersten $5\;\text{s}$?
\item Zwischen $5\;\text{s}$ und $12\;\text{s}$ wirkt ein Fahrwiderstand von $500\;\text{N}$. Wie gross ist dort die Gesamtkraft, und wie gross die Antriebskraft? Begründe.
\item Wie gross ist die Gesamtkraft beim Bremsen (mit Vorzeichen, Fahrtrichtung positiv)?
\end{enumerate}
\lpplatz{9}
\end{lpaufgabe}

\begin{lpaufgabe}{G3}{Schlitten}{3 P}
Ein Schlitten ($40\;\text{kg}$) wird aus dem Stand $4\;\text{s}$ lang mit $150\;\text{N}$ gezogen; die Reibung bremst mit $70\;\text{N}$.
Wie schnell ist er nach $4\;\text{s}$? Dann wird das Seil losgelassen, die Reibung bleibt. Wie lange gleitet er noch bis zum Stillstand?
\lpplatz{8}
\end{lpaufgabe}

\lpteil{Teil B · Gewichtskraft und Aufzug}{Kapitel 3 · 4 Punkte}

\begin{lpaufgabe}{G4}{Waage im Aufzug}{4 P}
Im ruhenden Aufzug zeigt eine Personenwaage $60\;\text{kg}$ an. Während der Fahrt zeigt sie einen Moment lang $52\;\text{kg}$.
Wie gross ist die Beschleunigung des Aufzugs in diesem Moment, und wohin zeigt sie?
Nenne zwei Fahrsituationen, in denen das vorkommt.
\lpplatz{9}
\end{lpaufgabe}

\lpteil{Teil C · Zwei Körper, ein Faden}{Kapitel 4 · 4 Punkte}

\begin{lpaufgabe}{G5}{Wagen und hängender Körper}{4 P}
Ein Wagen ($m_1 = 2.5\;\text{kg}$) steht auf einem reibungsfreien Tisch. Über eine Rolle zieht ihn ein Körper, der am Faden hängt.
Gemessen wird die Beschleunigung $a = 2.8\;\text{m/s}^2$. Welche Masse hat der hängende Körper? Wie gross ist die Fadenkraft?
Warum ist sie kleiner als die Gewichtskraft des hängenden Körpers?
\lpplatz{10}
\end{lpaufgabe}

\lpteil{Teil D · Kreisbewegung}{Kapitel 5 · 5 Punkte}

\begin{lpaufgabe}{G6}{Auto in der Kurve}{5 P}
Ein Auto ($1100\;\text{kg}$) fährt durch eine flache Kurve mit $r = 40\;\text{m}$. Die Haftreibung zwischen Reifen und Strasse kann höchstens $8800\;\text{N}$ betragen.
Wie schnell darf das Auto höchstens fahren (in $\text{m/s}$ und $\text{km/h}$)? Welche Rolle spielt die Haftreibung, und wohin zeigt sie?
Was geschieht, wenn das Auto schneller fährt?
\lpplatz{12}
\end{lpaufgabe}

\vfill
{\small\color{lpgrau}Bewerten: mit dem Bewertungspaket (Musterlösung, Punkteraster, Auftrag an eine KI) — erst nach dem Lösen öffnen.}
\end{document}
